\clearpage
\begin{thebibliography}{99}
\setlength{\itemsep}{5pt}

\bibitem{ridge}
Hoerl, A. E., \& Kennard, R. W. (1970).
Ridge regression: Biased estimation for nonorthogonal problems.
\emph{Technometrics}, 12(1), 55--67.
\href{https://doi.org/10.1080/00401706.1970.10488634}{DOI: 10.1080/00401706.1970.10488634}.
用途：Ridge 正則化背景；本文第 3 節依本案損失尺度給出矩陣推導。

\bibitem{gam}
Hastie, T., \& Tibshirani, R. (1986).
Generalized additive models.
\emph{Statistical Science}, 1(3), 297--310.
\href{https://doi.org/10.1214/ss/1177013604}{DOI: 10.1214/ss/1177013604}.
用途：加法模型的基本形式；不代表本文所有基底與門檻皆取自原文。

\bibitem{pspline}
Eilers, P. H. C., \& Marx, B. D. (1996).
Flexible smoothing with B-splines and penalties.
\emph{Statistical Science}, 11(2), 89--102（正文；討論續至 121 頁）.
\href{https://doi.org/10.1214/ss/1038425655}{DOI: 10.1214/ss/1038425655}.
用途：B-spline 係數差分懲罰；本文另加入全係數 L2 穩定項。

\bibitem{xgb}
Chen, T., \& Guestrin, C. (2016).
XGBoost: A scalable tree boosting system.
\emph{Proceedings of KDD}, 785--794.
\href{https://arxiv.org/abs/1603.02754}{作者論文：arXiv:1603.02754}.
用途：第 2.2 節正則化目標、二階展開、最佳葉值及分裂增益；
本文將其寫成加權平方損失的特例。

\bibitem{breiman}
Breiman, L. (1996).
Stacked regressions.
\emph{Machine Learning}, 24, 49--64.
\href{https://doi.org/10.1007/BF00117832}{DOI: 10.1007/BF00117832}.
用途：以交叉驗證預測學習非負線性組合的經典依據；
不保證任意資料都優於最佳單模型。

\bibitem{pystacked}
Ahrens, A., Hansen, C. B., \& Schaffer, M. E. (2023).
pystacked: Stacking generalization and machine learning in Stata.
\emph{The Stata Journal}, 23(4), 909--931.
\href{https://doi.org/10.1177/1536867X231212426}{DOI: 10.1177/1536867X231212426}.
用途：第 2 節 stacking 與限制方式；包含非負但不強制總和為 1 的選項。

\bibitem{bertsekas}
Bertsekas, D. P. (2009).
\emph{Convex Optimization Theory}. Athena Scientific.
\href{https://web.mit.edu/dimitrib/www/convexduality.html}{作者教材與相關資源（MIT）}.
用途：有限維凸分析、有限生成錐的封閉性；
本文附錄 A 提供本案所需的完整獨立證明。

\bibitem{boyd}
Boyd, S., \& Vandenberghe, L. (2004).
\emph{Convex Optimization}. Cambridge University Press.
\href{https://web.stanford.edu/~boyd/cvxbook/}{作者開放教材（Stanford）}.
用途：第 4--5 章凸最佳化與 KKT；
本文明確另證可行性、存在性，未僅由凸性推論有解。

\bibitem{fpp}
Hyndman, R. J., \& Athanasopoulos, G. (2021).
\emph{Forecasting: Principles and Practice}, 3rd ed. OTexts.
\href{https://otexts.com/fpp3/tscv.html}{第 5.10 節：時間序列交叉驗證}.
用途：只使用過去資料的向前驗證；活動分群與年度切分為本案設計。

\bibitem{cawley}
Cawley, G. C., \& Talbot, N. L. C. (2010).
On over-fitting in model selection and subsequent selection bias in performance evaluation.
\emph{Journal of Machine Learning Research}, 11, 2079--2107.
\href{https://www.jmlr.org/papers/v11/cawley10a.html}{JMLR 原文}.
用途：模型選擇與反覆評估可能造成的樂觀偏差。

\bibitem{little}
Little, J. D. C. (1961).
A proof for the queuing formula: $L=\lambda W$.
\emph{Operations Research}, 9(3), 383--387.
\href{https://doi.org/10.1287/opre.9.3.383}{DOI: 10.1287/opre.9.3.383}.
用途：長期排隊量、到達率與時間的關係；
本文第 7 節另證有限視野面積恆等式，未直接套用穩態假設。

\bibitem{urban}
Larson, R. C., \& Odoni, A. R. (1981).
\emph{Urban Operations Research}. Prentice-Hall.
\href{https://web.mit.edu/urban_or_book/www/book/chapter4/4.4.html}{MIT 開放教材，第 4.4 節}.
用途：排隊系統計數與 Little 關係的面積解釋。

\Needspace{6\baselineskip}
\bibitem{pan}
Pan, H., Liu, Y., Liu, Z., \& Hu, H. (2022).
Passenger inflow control along rail corridor for recurrent heavy passenger flow mitigation.
\emph{Transportation Research Record}, 2676(9), 59--75.
\href{https://doi.org/10.1177/03611981221085527}{DOI: 10.1177/03611981221085527}.
用途：共同考慮等待與營運成本、MIP 與 MPC 的研究方向；
本次核對摘要，本文式 \eqref{eq:dispatch} 是自行設計，非原文公式重現。

\bibitem{mol}
勞動部（2026）。基本工資／最低工資之制訂與調整經過。
\href{https://www.mol.gov.tw/1607/28690/2282/2284/2290/7121/}{勞動部官方資料}.
用途：2026 年 1 月 1 日起最低時薪 196 元；
換算乘客時間價值是本案研究假設，非該頁提出的交通估價方法。

\bibitem{lasso}
Tibshirani, R. (1996).
Regression shrinkage and selection via the lasso.
\emph{Journal of the Royal Statistical Society: Series B}, 58(1), 267--288.
\href{https://doi.org/10.1111/j.2517-6161.1996.tb02080.x}{DOI: 10.1111/j.2517-6161.1996.tb02080.x}.
用途：附錄 B 的 L1 懲罰背景；存在性論證在本文直接給出。

\bibitem{aic}
Akaike, H. (1974).
A new look at the statistical model identification.
\emph{IEEE Transactions on Automatic Control}, 19(6), 716--723.
\href{https://doi.org/10.1109/TAC.1974.1100705}{DOI: 10.1109/TAC.1974.1100705}.
用途：附錄 B 的 AIC 形式與資訊準則背景。
\end{thebibliography}
